% ============================================================================
%  Kernel V4 — Bilinear Projection Form for Radial Density
%  VSEPR-SIM Theoretical Division
%  Classification: TOP SECRET (internal development document)
%
%  Subject: Corrected separability analysis, gradient/numerical form,
%           matrix presolve form, and generalised bilinear projection
%           engine for multi-channel radial density computation.
%
%  Compile: pdflatex (twice for ToC)
% ============================================================================
\documentclass[11pt, a4paper, oneside]{article}

\usepackage[top=1.0in, bottom=1.0in, left=1.15in, right=1.15in,
			headheight=14pt]{geometry}
\usepackage[T1]{fontenc}
\usepackage{lmodern}
\usepackage{microtype}
\usepackage{amsmath, amssymb, amsthm, mathtools}
\usepackage{bm}
\usepackage{booktabs}
\usepackage{array}
\usepackage{tabularx}
\usepackage{listings}
\usepackage{xcolor}
\usepackage{enumitem}
\usepackage{fancyhdr}
\usepackage{titlesec}
\usepackage[colorlinks=true, linkcolor=black, citecolor=black,
			urlcolor=blue]{hyperref}

% --- Code style --------------------------------------------------------------
\definecolor{codegray}{rgb}{0.5,0.5,0.5}
\definecolor{codeblue}{rgb}{0.13,0.13,0.9}
\definecolor{codegreen}{rgb}{0.0,0.5,0.0}
\definecolor{codebg}{rgb}{0.97,0.97,0.97}
\lstset{
  backgroundcolor=\color{codebg},
  basicstyle=\ttfamily\small,
  keywordstyle=\color{codeblue}\bfseries,
  commentstyle=\color{codegreen}\itshape,
  numberstyle=\tiny\color{codegray},
  numbers=left, numbersep=6pt,
  breaklines=true, frame=single,
  rulecolor=\color{codegray},
  captionpos=b, showstringspaces=false,
  tabsize=2
}

% --- Headers -----------------------------------------------------------------
\pagestyle{fancy}
\fancyhf{}
\fancyhead[L]{\small\textit{Kernel V4 — Bilinear Projection Form}}
\fancyhead[R]{\small\textit{VSEPR-SIM / TOP SECRET}}
\fancyfoot[C]{\thepage}
\renewcommand{\headrulewidth}{0.4pt}

% --- Section style -----------------------------------------------------------
\titleformat{\section}{\Large\bfseries}{\S\thesection}{1em}{}
\titleformat{\subsection}{\large\bfseries}{\thesubsection}{1em}{}
\titleformat{\subsubsection}{\normalsize\bfseries}{\thesubsubsection}{1em}{}

% --- Math convenience --------------------------------------------------------
\newcommand{\Ri}[1]{R_{#1}^2(\mathbf{r})}          % R_k^2(r)
\newcommand{\Rit}[2]{R_{#1}^2(\mathbf{r}_{#2})}     % R_k^2(r_i)
\newcommand{\ez}{e_z(t)}
\newcommand{\Bvec}{B(\mathbf{r})}
\newcommand{\Dvec}{D(\mathbf{r})}
\newcommand{\wvec}{\mathbf{w}(t)}
\newcommand{\Rvec}{\mathbf{R}(\mathbf{r})}
\newcommand{\Rgrid}{R_{\mathrm{grid}}}
\newcommand{\Mmat}{\mathbf{M}}
\newcommand{\evec}{\mathbf{e}(t)}

% ============================================================================
\begin{document}

\begin{titlepage}
  \centering
  \vspace*{1.2cm}
  {\large\scshape VSEPR-SIM Theoretical Division\par}
  \vspace{0.3cm}
  {\small TOP SECRET — Internal Development Document\par}
  \vspace{1.2cm}
  {\Huge\bfseries Kernel V4\par}
  \vspace{0.4cm}
  {\LARGE Bilinear Projection Form\\[0.2em]
   for Radial Density\par}
  \vspace{1.2cm}
  {\normalsize
	Corrected separability analysis\,|\,
	Gradient/numerical caching form\,|\,\\[0.3em]
	Matrix presolve form\,|\,
	Generalised multi-channel projection engine\par}
  \vspace{1.5cm}
  {\normalsize VSEPR-SIM Development Day 64\par}
\end{titlepage}

\tableofcontents
\newpage

% ============================================================================
\section{Correction: Separability in the Gradient Form}
\label{sec:correction}
% ============================================================================

\subsection{The claim that needed fixing}

An earlier version of the derivation stated that the gradient form
\emph{destroys} separability of the density into spatial and temporal
parts, and therefore cannot cache spatial structure. This is too strong.

The gradient form hides the clean matrix structure, but separability
is preserved. Define the \textbf{base density} and \textbf{difference density}:

\begin{align}
  \Bvec &= \tfrac{1}{2}\bigl[\Ri{1} + \Ri{2}\bigr]
  \label{eq:B} \\[4pt]
  \Dvec &= \tfrac{1}{2}\bigl[\Ri{1} - \Ri{2}\bigr]
  \label{eq:D}
\end{align}

Then the full density is:

\begin{equation}
  \boxed{\rho(\mathbf{r},t) = \Bvec + \ez\,\Dvec}
  \label{eq:rho_BD}
\end{equation}

This is already cacheable. Spatial derivatives separate cleanly:

\begin{equation}
  \nabla\rho(\mathbf{r},t) = \nabla\Bvec + \ez\,\nabla\Dvec
  \label{eq:grad_rho}
\end{equation}

since $\ez$ carries no spatial dependence.

\medskip
\noindent\textbf{Corrected claim:}
\begin{quote}
  \itshape
  The gradient form does not destroy separability; it merely hides the
  clean matrix structure. The matrix form makes the separability
  \emph{explicit} and \emph{generalisable}. The two claims are
  complementary, not equivalent.
\end{quote}

% ============================================================================
\section{Three Equivalent Forms}
\label{sec:three_forms}
% ============================================================================

\subsection{Form 1 — Decomposed projection form (readable)}

\begin{equation}
  \rho(\mathbf{r},t)
  = \tfrac{1}{2}\bigl[(1+\ez)\,\Ri{1} + (1-\ez)\,\Ri{2}\bigr]
  \label{eq:form1}
\end{equation}

This is the most readable form. It shows directly that the two radial
projection channels $R_1^2$ and $R_2^2$ are weighted by
$(1\pm\ez)/2$, which sum to 1 when $\ez = 0$ (equal weight)
and collapse to a single channel when $\ez = \pm 1$.

\subsection{Form 2 — Expanded modulation form (numerical / caching)}

Expand \eqref{eq:form1}:

\begin{equation}
  \rho(\mathbf{r},t)
  = \underbrace{\tfrac{1}{2}\Ri{1} + \tfrac{1}{2}\Ri{2}}_{\Bvec}
  + \ez\underbrace{\tfrac{1}{2}\bigl[\Ri{1}-\Ri{2}\bigr]}_{\Dvec}
  \label{eq:form2_expand}
\end{equation}

giving \eqref{eq:rho_BD}.

\medskip
\noindent\textbf{Precomputed cache per grid point} $\mathbf{r}_i$:

\begin{equation}
  \bigl\{B_i,\; D_i,\; \nabla B_i,\; \nabla D_i\bigr\}
  \label{eq:cache}
\end{equation}

\noindent\textbf{Runtime evaluation at time} $t$:

\begin{align}
  \rho_i(t)       &= B_i + \ez\,D_i          \label{eq:rt_rho} \\
  \nabla\rho_i(t) &= \nabla B_i + \ez\,\nabla D_i  \label{eq:rt_grad}
\end{align}

This is the computationally useful statement: spatial precomputation is
done once; each timestep costs two flops per grid point (one multiply,
one add) for $\rho$, and $2d$ flops for $\nabla\rho$ in $d$ spatial
dimensions.

\subsection{Form 3 — Matrix presolve form (dot-product runtime)}

Write the radial channel vector and time-weight vector:

\begin{equation}
  \Rvec =
  \begin{bmatrix} \Ri{1} \\ \Ri{2} \end{bmatrix},
  \qquad
  \wvec = \tfrac{1}{2}
  \begin{bmatrix} 1+\ez \\ 1-\ez \end{bmatrix}
  \label{eq:Rw}
\end{equation}

Then:

\begin{equation}
  \boxed{\rho(\mathbf{r},t) = \Rvec^\top \wvec}
  \label{eq:form3}
\end{equation}

Runtime cost for a single preloaded grid point: $O(1)$.

% ============================================================================
\section{Generalised Bilinear Projection Form}
\label{sec:general}
% ============================================================================

For $n$ radial channels and $m+1$ temporal state channels, define:

\begin{equation}
  \Rvec =
  \begin{bmatrix} \Ri{1} \\ \Ri{2} \\ \vdots \\ R_n^2(\mathbf{r}) \end{bmatrix}
  \in \mathbb{R}^n,
  \qquad
  \evec =
  \begin{bmatrix} 1 \\ e_1(t) \\ e_2(t) \\ \vdots \\ e_m(t) \end{bmatrix}
  \in \mathbb{R}^{m+1}
  \label{eq:Re_general}
\end{equation}

Let $\Mmat \in \mathbb{R}^{n \times (m+1)}$ be the \textbf{channel mixing matrix}
that maps temporal state channels into radial-channel weights. Then:

\begin{equation}
  \boxed{\rho(\mathbf{r},t) = \Rvec^\top \Mmat\, \evec(t)}
  \label{eq:bilinear}
\end{equation}

\medskip
The two-channel case \eqref{eq:form3} is recovered with:
\begin{equation}
  \Mmat = \tfrac{1}{2}
  \begin{bmatrix}
	1 &  1 \\
	1 & -1
  \end{bmatrix},
  \qquad
  \evec(t) =
  \begin{bmatrix} 1 \\ \ez \end{bmatrix}
  \label{eq:M2}
\end{equation}

\medskip
\noindent\textbf{Why this form matters for code architecture:}
\begin{itemize}
  \item $\Rvec$ and $\Mmat$ are fully spatial / static — precomputed
	once before the time loop.
  \item $\evec(t)$ is a small vector updated each step from
	the temporal state kernel.
  \item The inner product $\Rvec^\top \Mmat\,\evec$ at each grid point
	is a dense matrix-vector product on a small system
	($n \times (m+1)$ is typically $2$--$8 \times 2$--$8$).
  \item This scales cleanly from a 2-channel toy case
	to a multi-channel projection engine without restructuring.
\end{itemize}

% ============================================================================
\section{Presolve Pipeline}
\label{sec:pipeline}
% ============================================================================

\subsection{Offline presolve (done once)}

For each spatial grid point $\mathbf{r}_i$, $i = 1, \ldots, N$:

\begin{equation}
  \mathbf{R}_i =
  \begin{bmatrix} \Rit{1}{i} \\ \Rit{2}{i} \end{bmatrix}
\end{equation}

Assemble the grid matrix:

\begin{equation}
  \Rgrid \in \mathbb{R}^{N \times n},
  \quad
  [\Rgrid]_{i,k} = R_k^2(\mathbf{r}_i)
\end{equation}

Compute and cache $\Rgrid \Mmat$ (static, shape $N \times (m+1)$).

\subsection{Runtime (each timestep)}

Update $\evec(t)$ from the temporal state kernel, then:

\begin{equation}
  \bm{\rho}(t) = (\Rgrid \Mmat)\,\evec(t)
  \label{eq:runtime_vec}
\end{equation}

\noindent\textbf{Cost analysis:}

\begin{center}
\begin{tabular}{lll}
  \toprule
  Operation & Cost & Notes \\
  \midrule
  Precompute $\Rgrid \Mmat$  & $O(N n (m+1))$ & Once, offline \\
  Update $\evec(t)$          & $O(m)$         & Temporal kernel \\
  Evaluate $\bm{\rho}(t)$    & $O(N(m+1))$    & Matrix-vector product \\
  Single-point lookup        & $O(m+1) \approx O(1)$ & After presolve \\
  \bottomrule
\end{tabular}
\end{center}

For the two-channel case $n = 2$, $m = 1$:
runtime evaluation is $O(2N)$, or $O(1)$ per preloaded point.

% ============================================================================
\section{Reference Implementation}
\label{sec:impl}
% ============================================================================

\subsection{Python (prototype)}

\begin{lstlisting}[language=Python, caption={Kernel V4 two-channel presolve (Python)}]
import numpy as np

class KernelV4:
	"""
	Bilinear projection engine for radial density.
	Two-channel case: rho(r,t) = R(r)^T w(t)
	"""

	def __init__(self, R1_grid: np.ndarray, R2_grid: np.ndarray):
		"""
		Parameters
		----------
		R1_grid : (N,) array of R1^2 values on the spatial grid
		R2_grid : (N,) array of R2^2 values on the spatial grid
		"""
		# Spatial cache  [N x 2]
		self.R_grid = np.stack([R1_grid, R2_grid], axis=1)

		# Channel mixing matrix M  [2 x 2]
		self.M = 0.5 * np.array([[1.0,  1.0],
								  [1.0, -1.0]])

		# Presolve:  R_grid @ M  is [N x 2], computed once
		self._RM = self.R_grid @ self.M

	def weight_vector(self, ez: float) -> np.ndarray:
		"""Temporal weight vector e(t) = [1, ez]."""
		return np.array([1.0, ez])

	def density(self, ez: float) -> np.ndarray:
		"""Runtime evaluation: rho(t) = (R_grid M) e(t).  Cost: O(2N)."""
		return self._RM @ self.weight_vector(ez)

	def density_single(self, i: int, ez: float) -> float:
		"""Single-point lookup after presolve.  Cost: O(1)."""
		return float(self._RM[i] @ self.weight_vector(ez))


# --- Example usage -----------------------------------------------------------
if __name__ == "__main__":
	N = 1000
	r = np.linspace(0.0, 5.0, N)

	# Toy radial functions: hydrogen-like s-orbitals
	R1 = np.exp(-r)             # R1^2 ~ exp(-2r), but use R1 for clarity
	R2 = r * np.exp(-r / 2.0)

	kernel = KernelV4(R1**2, R2**2)

	# Time evolution: ez sweeps from -1 to 1
	for ez in np.linspace(-1.0, 1.0, 5):
		rho = kernel.density(ez)
		print(f"ez={ez:+.2f}  rho[0]={rho[0]:.6f}  rho[-1]={rho[-1]:.6f}")
\end{lstlisting}

\subsection{C++ kernel hook (production sketch)}

\begin{lstlisting}[language=C++, caption={Kernel V4 presolve hook (C++23)}]
#pragma once
#include <span>
#include <vector>
#include <array>
#include <cstddef>

namespace ikk {

/// Two-channel bilinear projection kernel.
/// rho(r,t) = R_grid * M * e(t)
/// Presolve: cache R_grid * M offline.  Runtime: O(2N) matvec.
class KernelV4 {
public:
	/// Presolve constructor.
	/// r1sq, r2sq: R1^2 and R2^2 evaluated on the spatial grid, length N.
	KernelV4(std::span<const double> r1sq,
			 std::span<const double> r2sq)
	{
		const std::size_t N = r1sq.size();
		rm_.resize(N);                  // rm_[i] = {B_i, D_i}

		for (std::size_t i = 0; i < N; ++i) {
			rm_[i][0] = 0.5 * (r1sq[i] + r2sq[i]);   // B_i
			rm_[i][1] = 0.5 * (r1sq[i] - r2sq[i]);   // D_i
		}
	}

	/// Full grid evaluation.  rho_out must have length N.
	/// Cost: O(2N)
	void density(double ez, std::span<double> rho_out) const noexcept {
		for (std::size_t i = 0; i < rm_.size(); ++i)
			rho_out[i] = rm_[i][0] + ez * rm_[i][1];
	}

	/// Single-point lookup after presolve.  Cost: O(1).
	[[nodiscard]]
	double density_at(std::size_t i, double ez) const noexcept {
		return rm_[i][0] + ez * rm_[i][1];
	}

private:
	/// Cached (R_grid * M): rm_[i] = {B_i, D_i}
	std::vector<std::array<double, 2>> rm_;
};

} // namespace ikk
\end{lstlisting}

% ============================================================================
\section{Summary}
\label{sec:summary}
% ============================================================================

\begin{center}
\begin{tabular}{llll}
  \toprule
  Form & Equation & Best for & Cache \\
  \midrule
  Decomposed projection  & \eqref{eq:form1}   & Reading, intuition        & No \\
  Expanded modulation    & \eqref{eq:rho_BD}  & Gradient / finite-diff    & $B_i, D_i, \nabla B_i, \nabla D_i$ \\
  Matrix presolve        & \eqref{eq:form3}   & Dot-product runtime       & $\mathbf{R}_i^\top$ \\
  Generalised bilinear   & \eqref{eq:bilinear}& Multi-channel engine      & $\Rgrid \Mmat$ \\
  \bottomrule
\end{tabular}
\end{center}

\vspace{1em}
\noindent
All four forms are algebraically equivalent for the two-channel case.
The generalised bilinear form \eqref{eq:bilinear} is the target for
production code because it is:

\begin{itemize}
  \item \textbf{Statically typed by shape}: $\Rgrid \in \mathbb{R}^{N \times n}$,
	$\Mmat \in \mathbb{R}^{n \times (m+1)}$, $\evec \in \mathbb{R}^{m+1}$.
  \item \textbf{BLAS-friendly}: the presolve is a single DGEMM;
	the runtime step is a DGEMV.
  \item \textbf{Channel-count agnostic}: adding a new radial or temporal
	channel is a shape change in $\Rgrid$ or $\Mmat$,
	not a restructuring of the kernel.
\end{itemize}

\vfill
\begin{center}
  \rule{0.4\textwidth}{0.4pt}\\[0.5em]
  \textit{Kernel V4 — VSEPR-SIM Theoretical Division.
  Humanity survives another algebraic rearrangement.}
\end{center}

\end{document}
